\documentclass[10pt,a4paper]{amsart}
\usepackage[margin=19mm]{geometry}
\usepackage{amsmath,amssymb,mathtools}
\usepackage[scr=rsfs,cal=euler]{mathalfa}
\usepackage{xfrac}
\usepackage[unicode,hidelinks,bookmarks=false]{hyperref}
\newcommand{\ie}{i.e.\ }
\newcommand{\cf}{cf.\ }
\newcommand{\resp}{respectively\ }
\newcommand{\Subsection}{\subsection}
\providecommand{\nobreakdash}{\nobreak\hbox{-}}
\setlength{\emergencystretch}{2em}
\multlinegap0pt
\usepackage{amssymb}
\usepackage{enumerate}
\newtheorem{theorem}{Theorem}
\newtheorem{lemma}{Lemma}
\title{Chainmail links and $L$-spaces}
\author{Ian Agol}
\date{Source: arXiv:2306.10918v1 (CC BY 4.0)}
\begin{document}
\maketitle

\noindent\emph{Source and licence.} Chainmail links and $L$-spaces. Version arXiv:2306.10918v1 (CC BY 4.0), distributed by arXiv under the Creative Commons Attribution 4.0 International licence, http://creativecommons.org/licenses/by/4.0/, as recorded in the arXiv licence metadata for this exact version and shown on the abstract page https://arxiv.org/abs/2306.10918v1. Full licence notice: \texttt{licenses/arxiv-2306.10918-CC-BY-4.0.txt}.\par\smallskip

\noindent\textit{Editorial scope.} Counted unit: the article's lemma lemma (source0.tex lines 346--349). The article's complete original proof of it is delivered with it (source0.tex lines 350--370, one \texttt{proof} environment of 21 lines). No mathematics is written, repaired or re-derived: the statement and the proof are the article's own lines, in the article's own notation, and no retained line is substituted or re-flowed.  Retained uncounted as the source-local context the proof needs: lines 248--252; lines 260--262.  One declared adaptation: exactly 1 ancillary strings inside the retained spans are omitted (image-inclusion commands and, where the source repeats a label, the repeated label definitions) and no other line differs from the source; the figure environments, with their placement, captions and labels, are kept, and the package records the omitted strings in fidelity receipts bound to the affected bodies.  No retained line contains a citation, so the wrapper carries no bibliography and no bibliography entry.  The retained lines contain the article's own diagram code, which the wrapper typesets with the standard tikz packages; no image file is delivered and the package declares no asset dependency.  Editorial formatting only, in addition: an AMS \texttt{amsart} preamble that restates the article's own declarations and macros used by the retained lines, namely the environments \texttt{theorem}, \texttt{lemma} on separate counters, together with the standard packages those lines need.  The retained environments print in this document as Theorem 1, Lemma 1 in order of appearance, whereas the article prints the same items with the numbering of its own full document.  The complete native source of the article as distributed in the arXiv e-print for this exact version is bundled unchanged as \texttt{sources/143084/source0.tex}.
\begin{figure}[htb]

\caption{Links and surgeries on augmented link corresponding to minors of $G$} \label{crossing_loop}

\end{figure}

\begin{theorem} \label{alternating}
Negative alternating chainmail links are $L$-space links. More specifically, for a chainmail graph $G$ with all $\nu(v)\geq 0, v\in V(G)$, and at least one $\nu(v)>0$ in each component of $G$, and $\epsilon(e) < 0$ for all $e\in E(G)$,  $M_G$ is an $L$-space and an instanton $L$-space.
\end{theorem}

\begin{lemma}
Consider a partially augmented chainmail link simplicial graph $G$, with negative weights on the vertices $V_c(G) \subset V(G)$ corresponding to crossing loops and non-negative weights on the other vertices $V(G)-V_c(G)$. Moreover, the edges between vertices of $V(G)-V_c(G)$ have weight $-1$. Then the sign of the determinant is $(-1)^{|V_c(G)|}$ if non-zero. 

\end{lemma}

\begin{proof}
The proof is by induction on the sum of the negative crossing loop weights. If there are no crossing loops, then $\Lambda(G)$ is is the linking matrix of a negative alternating chainmail graph, and hence the sign of the determinant is non-negative by the proof of Theorem \ref{alternating}, establishing the base case. 

Now consider a single crossing loop of the flat augmented link with weight $- c$, $c > 0$, and its two adjacent loops having framings $a'$ and  $b'$. The matrix $\Lambda(G)$ looks like 

$$\left(\begin{array}{cccc}-c & 1 & -1 & 0 \\1 & a' & 0 & \ast \\-1 & 0 & b' & \ast \\0 & \ast & \ast & A\end{array}\right).$$

When $c=1$, we do column operations adding the first column to the second and subtracting the first column from the third to get 

 $$\left|\begin{array}{cccc}-1 & 0 & 0 & 0 \\1 & a'+1 & -1 & \ast \\-1 & -1 & b'+1 & \ast \\0 & \ast & \ast & A\end{array}\right| = (-1)\left|\begin{array}{ccc}a'+1 & -1 & \ast \\ -1 & b'+1 & \ast \\ \ast & \ast & A\end{array}\right|  .$$

The last determinant is that of a partially augmented chainmail link by doing a $-1$ twist along the crossing loop (see the left diagrams of Figure \ref{crossing_loop}). So the sign of this determinant is $(-1)^{(|V_c(G)|-1)}$ by induction. Thus it is true by induction for $c=1$. 

When $c>1$, we have 

$$\left|\begin{array}{cccc}-c-1 & 1 & -1 & 0 \\1 & a' & 0 & \ast \\-1 & 0 & b' & \ast \\0 & \ast & \ast & A\end{array}\right| =\left|\begin{array}{cccc}-c & 1 & -1 & 0 \\1 & a' & 0 & \ast \\-1 & 0 & b' & \ast \\0 & \ast & \ast & A\end{array}\right| - \left|\begin{array}{cccc}1&0 & 0 & 0 \\1 & a' & 0 & \ast \\-1 & 0 & b' & \ast \\0 & \ast & \ast & A\end{array}\right| .$$
The matrix on the right has determinant 

$$\left|\begin{array}{ccc}  a' & 0 & \ast \\ 0 & b' & \ast \\ \ast & \ast & A\end{array}\right|$$
which is $det \Lambda(G-v)$. By induction, this matrix has the opposite sign, so we get by induction on $c$ that the sign with weight $-c-1$ is $(-1)^{|V_c(G)|}$. 
\end{proof}

\end{document}
